\chapter{Thalès et triangles semblables}
\label{corriges:14}
\section*{Exercice 14.1 -- Configuration emboîtée}
Thalès donne $AM/AB=10/15=2/3=AN/AC=MN/BC$. Ainsi $AN=(2/3)12=8\unite{cm}$ et $MN=(2/3)18=12\unite{cm}$. Comme $M$ appartient au segment $[AB]$, $MB=AB-AM=5\unite{cm}$.

\section*{Exercice 14.2 -- Prouver un parallélisme}
Les alignements et la condition de position sont vérifiés. Les quotients $7/10$ et $10{,}5/15$ valent tous deux $0{,}7$. La réciproque de Thalès donne le parallélisme. Si $AN=11$, $11/15\ne7/10$ ; le théorème direct serait contredit par un parallélisme. Les droites ne seraient donc pas parallèles.

\section*{Exercice 14.3 -- Configuration papillon}
Le rapport est $AM/AB=2/3$. Donc $AN=4$ et $MN=8$. Puisque $A$ est entre les extrémités, $BM=BA+AM=9+6=15$ et $CN=CA+AN=6+4=10$. Soustraire les longueurs dans ce papillon serait une erreur de position.

\section*{Exercice 14.4 -- Angles et similitude}
Les angles correspondent deux à deux, donc les triangles sont semblables. Le rapport de longueur est $9/6=3/2$. Le périmètre vaut $24\times3/2=36$. L'aire vaut $18\times(3/2)^2=18\times9/4=40{,}5$. On associe les côtés par les angles opposés, pas par leur orientation dans un dessin.
